\Askhsh{1790}{4}
\begin{ylh}{1}
    \item 3.2. 1ο Κριτήριο ισότητας τριγώνων
    \item 4.2. Τέμνουσα δύο ευθειών - Ευκλείδειο αίτημα
    \item 5.2. Παραλληλόγραμμα
    \item 5.6. Εφαρμογές στα τρίγωνα
    \item 5.10. Τραπέζιο
    \item 5.11. Ισοσκελές τραπέζιο.
\end{ylh}
Δίνεται παραλληλόγραμμο $\mathrm{A}\mathrm{B}\Gamma\Delta$ με τη γωνία του $\hat{\mathrm{B}}$ να είναι ίση με $70^\circ$ και το ύψος του $\mathrm{A}\mathrm{E}$. Έστω $\mathrm{H}$ σημείο της $\mathrm{B}\Gamma$ ώστε $\mathrm{B}\mathrm{E} = \mathrm{E}\mathrm{H}$.
\begin{alist}
    \item Να αποδείξετε ότι το τετράπλευρο $\mathrm{A}\mathrm{H}\Gamma\Delta$ είναι ισοσκελές τραπέζιο. \monades{8}
    \item Να υπολογίσετε τις γωνίες του τραπεζίου $\mathrm{A}\mathrm{H}\Gamma\Delta$. \monades{9}
    \item Αν $\mathrm{M}$ το μέσο της $\mathrm{B}\Delta$, να αποδείξετε ότι $\mathrm{E}\mathrm{M} = \frac{\mathrm{A}\Gamma}{2}$. \monades{8}
\end{alist}
